\section{Introduction}

In December 2020, FireEye disclosed that a backdoor had been distributed to customers of SolarWinds inside a correctly signed update of the Orion IT-monitoring platform. The campaign was attributed to the SVR by U.S. and UK authorities, and MITRE tracks the activity as APT29~\cite{mitre-apt29,aa22-110a}. The incident matters beyond its victim list because nothing in the delivery path was broken. The supplier's build system produced the malicious component, the supplier's signing process signed it, and customers installed it through the update channel they had been told to trust.

The incident is reconstructed from public reporting, with each event tied to the source that reported it. The recreation in Section~\ref{sec:recreation} follows the behaviour Mandiant and CrowdStrike documented; it does not run the actor's binaries~\cite{mandiant-sunburst-details,crowdstrike-sunspot}.

The report covers the group's history, the techniques and infrastructure behind the incident, the victims and their losses, and a lab that reproduces the trust failure rather than the malware. It closes with six peer-reviewed papers and a research question. A valid signature answers whether the supplier signed an artefact. It does not answer whether the build that produced the artefact was clean.
